\documentclass[20pt]{article}
\usepackage[paperwidth=1189mm,paperheight=841mm,margin=24mm]{geometry}
\usepackage[T1]{fontenc}
\usepackage{lmodern}
\usepackage{xcolor}
\usepackage{tikz}
\usetikzlibrary{positioning}
\usepackage{microtype}
\pagestyle{empty}
\definecolor{accent}{HTML}{143D6B}
\definecolor{highlight}{HTML}{C93C4B}
\definecolor{panel}{HTML}{EDF4FC}
\definecolor{ink}{HTML}{182331}
\newcommand{\head}[1]{\vspace{2mm}{\Large\bfseries\color{accent}#1}\par\vspace{2mm}}
\begin{document}
\color{ink}
\begin{tikzpicture}[remember picture,overlay]
  \fill[accent] (current page.north west) rectangle ([yshift=-35mm]current page.north east);
  \fill[highlight] ([yshift=-35mm]current page.north west) rectangle ([yshift=-40mm]current page.north east);
\end{tikzpicture}
\vspace*{-7mm}
\color{white}{\fontsize{38}{42}\selectfont\bfseries Electrical and Electronic Engineering}\par
\vspace{5mm}{\fontsize{25}{30}\selectfont\bfseries A low-power sensing node for crowded environments}\par
\vspace{3mm}{\Large Taylor Morgan \quad\textbullet\quad Final-Year Project}
\vspace{14mm}\color{ink}
\noindent
\begin{minipage}[t]{0.31\linewidth}
\head{Context}
Research is easier to communicate when the question, method, and decision are visible together. This starter gives a technical project a clear narrative with enough space for evidence and limitations.
\head{Question}
How can a compact, repeatable workflow make results easier for another reader to evaluate?
\head{Method}
\colorbox{panel}{\parbox{0.94\linewidth}{\large\textbf{Question} $\rightarrow$ \textbf{Measure} $\rightarrow$ \textbf{Compare} $\rightarrow$ \textbf{Decide}}}
\head{Study design}
We compared a baseline process with the proposed workflow using repeated trials and a fixed review protocol. Replace this sample paragraph with the details of your own project.
\end{minipage}
\hfill
\begin{minipage}[t]{0.31\linewidth}
\head{Results}
\begin{tabular}{p{0.52\linewidth}rr}
Measure & Baseline & Proposed\\\hline
Agreement & 58\% & \textbf{88\%}\\
Review time & 26 min & \textbf{15 min}\\
Complete records & 63\% & \textbf{94\%}
\end{tabular}
\head{Main finding}
\colorbox{panel}{\parbox{0.94\linewidth}{\centering{\fontsize{42}{46}\selectfont\bfseries\color{highlight}2.1$\times$}\par\Large faster shared decisions}}
\head{Interpretation}
The largest gain came from making assumptions explicit at the start. The result remained stable across the repeated trials, while the review process became easier to audit.
\head{Evidence}
Use this space for a figure, a short table, or one carefully chosen visualization. Keeping the evidence close to its interpretation helps readers follow the argument.
\end{minipage}
\hfill
\begin{minipage}[t]{0.31\linewidth}
\head{Conclusion}
A good poster is a coordination tool: it makes the next action visible without hiding uncertainty. Replace this conclusion with the strongest defensible claim from your work.
\head{Limitations}
The sample is deliberately small and the test setting is controlled. Transfer to other populations, environments, or operating conditions should be tested before making a broader claim.
\head{Next steps}
\begin{itemize}\setlength\itemsep{3mm}
\item Extend the evaluation to new conditions.
\item Release the smallest reproducible artifact.
\item Compare independent teams or sites.
\end{itemize}
\head{Takeaway}
\colorbox{panel}{\parbox{0.94\linewidth}{\Large\bfseries\color{accent}Make the evidence trail easy to follow.}}
\head{Contact}
taylor.morgan@example.org\\
\texttt{example.org/project}
\end{minipage}
\vfill\hfill{\small EEE-style starter | Final-Year Project 2026}
\end{document}
